\section{Additional established strict sectors in the same arithmetic window}
\label{sec:additional-established-window-sectors}

\paragraph{Scope.}
This addendum collects the additional families established by the
specific earlier inputs audited below.  It is not an assertion that
every possible consequence of every installed theorem has been exhausted.
No new class computation, regulator estimate, or coefficient-field
Bloch--Kato implication is used.
The character labels, coefficient rings, full local conditions, and
common window remain exactly those of AC-K25 and AC-K35.
The proof is a single extension of AC-G145-window; its frozen proof
and the strict theorems are not modified.

Use $p=269$, $K=\mathbb Q(\omega)$, $E_{\rm c}=K(\sqrt[3]{269})$,
and its cubic character $\chi$.
Let $\phi$ be the conductor-sixteen character of AC-G142/AC-G144,
and put $\psi_4=\chi_{-4}|_{G_K}$.
Use the actual host $H_u$, group $G_u$, rings $O_\ell$, fraction
fields $E_\ell$, degrees $b_\ell$, and $T=\lfloor u^{399}\rfloor$
of AC-G145.  Take $u\geq2$, so $h(u)\geq2$.
For any positive integer $T$ define the following disjoint label sets:
\[
 \begin{aligned}
 Z_{4,16}={}&\{(p,\chi^{-j}\theta,e):
       j=1,2,\ \theta\in\{\psi_4,\phi,\phi^{-1}\},\\
       &\hspace{35mm}1\leq e<T,\ e\equiv-1\pmod{134}\},\\
 Z_{\rm cub}={}&\{(p,\chi^j,e):
       j=0,1,2,\ 1\leq e<T,\ e\equiv0\text{ or }267\pmod{268}\},\\
 Z_3={}&\{(3,1,e):1\leq e<T\},\\
 Z_{\rm est}={}&Z_{4,16}\ \sqcup\ Z_{\rm cub}\ \sqcup\ Z_3 .
 \end{aligned}
 \tag{EW.1}
\]
The counts in the respective rows are
\[
 \begin{aligned}
 |Z_{4,16}|&=6\lfloor T/134\rfloor,\\
 |Z_{\rm cub}|&=
       3\bigl(\lfloor(T-1)/268\rfloor+\lfloor T/268\rfloor\bigr),\\
 |Z_3|&=T-1 .
 \end{aligned}
 \tag{EW.2}
\]
Let $I_{399}^{\rm est}$ omit precisely $Z_{\rm est}$ from (TC.1).
Every unlisted coefficient prime, character and weight remains.

\begin{corollary}[AC-G146: additional established sectors in the actual window]
\label{AC-G146}
Every label in (EW.1) occurs on the actual prescribed host, and
its strict group (FC.2) is zero over the original $O_\ell$.
The exact counts are (EW.2).
Writing $D=|G_u|$ and retaining (FL.1), one has the single bound
\[
 \begin{split}
 0\leq I_{399}-I_{399}^{\rm est}
 \leq{}&8D^2T^2\log(8TD)\\
 &+\tfrac12D\mathcal R T^2+2Dr c_0T\\
 ={}&O_\varepsilon\bigl(u^{802+16\varepsilon}(\log u)^2\bigr).
 \end{split}
 \tag{EW.3}
\]
Thus the identical common-window premise with $I_{399}^{\rm est}$
is sufficient for (TC.3), with half the relative slack and only an
increased initial threshold.  The converse is immediate.
The same statement holds relative to either previously reduced
conductor-sixteen sum.
No estimate for the remaining sum or the mixed realization is asserted.
\end{corollary}

\begin{proof}
\proofstep{1. Actual characters and the original strict kernel.}
The fields $\mathbb Q(i)$ and $\mathbb Q(\zeta_{16})$ are contained
in the two-power cyclotomic factor of $H_u$.
The field $E_{\rm c}$ is its actual cubic factor.
Consequently every finite character in (EW.1) factors through $G_u$.
The orders of $\psi_4$ and $\phi^{\pm1}$ are respectively two and four,
and their restrictions are nontrivial over $K$.
They remain distinct from the characters trivial on the noncubic
factor.  The cubic subgroup separates the nontrivial cubic labels.
The coefficient $3$ row is disjoint from every coefficient $269$ row.
This proves the claimed set disjointness.
Both coefficient primes belong to the unrestricted prime range of (TC.1).
The exponent of $G_u$ contains three and four, so its original
coefficient rings already contain the displayed character values.

Every coefficient is the original one
\[
 A_{\ell,\eta,e}
 =(O_\ell\otimes_{\mathbb Z_\ell}\mathbb Q_\ell/\mathbb Z_\ell)
                      (\eta)(-e)
 =(E_\ell/O_\ell)(\eta c_\ell^{-e}).
 \tag{EW.4}
\]
The scalar action on $O_\ell$ is trivial; its abstract roots of
unity are coefficient values, not a second Galois action.
No replacement by its residue field, unramified subring, or
degree-one coefficient ring is made.
At $\ell$ the strict condition is zero in full local $H^1$;
elsewhere it is zero on inertia.

We use the same enlargement comparison as AC-G143.
At an added away-$\ell$ place where (EW.4) is unramified,
a zero inertia cohomology class is a zero inertia cocycle.
It vanishes on the closed normal subgroup generated by such
inertia groups and descends to the smaller Galois group.
Conversely inflation is injective and preserves all local conditions.
Thus a global full-local strict-zero theorem, or the corresponding
minimal unramified-outside-$\ell$ statement, applies to the actual
$\Sigma_\ell$ in (FC.2).

\proofstep{2. The conductor-four and conductor-sixteen rows.}
For $\theta=\psi_4$, apply AC-G114-strict and AC-G115-strict,
or their exact two-residue version AC-G125-lifts with $r=0$ and
$\kappa=1$.
They state original-field strict vanishing for
$(E_p/O_p)(\chi^{-j}\psi_4 c_p^{-e})$
when $e\equiv133$ or $267\bmod268$.
They impose no coefficient ramification condition.
The two quartic rows are precisely AC-G142 and AC-G144, with
their additional finite exponent zero and continuous character one.
Their original coefficients and local conditions are (EW.4).
This proves all six rows of $Z_{4,16}$.
The two progressions are exactly $e=134n-1$,
giving the first count in (EW.2).

\proofstep{3. The old fixed-cubic-field branch at residue zero.}
AC-K31, equation (K31.2), proves
\[
             S_p(E_{\rm c},e)=0
                 \qquad(e\geq1,\ e\equiv0\bmod268).
 \tag{EW.5}
\]
Here its strict group is unramified outside $p$ and zero in full
cohomology at every $p$-place, with coefficient
$\mathbb Q_p/\mathbb Z_p(-e)$.
The local invariant correction needed for (EW.5) is part of the
complete AC-K31 proof, not an omitted local term.

For $j=0,1,2$, restrict a strict class with coefficient
$(E_p/O_p)(\chi^j c_p^{-e})$ from $K$ to $E_{\rm c}$.
The cubic character becomes trivial.
Every full local zero and inertia-zero condition remains zero
on the corresponding subgroups upstairs.
After removing any added unramified primes as in Step1, the class lies in
\[
                  O_p\otimes_{\mathbb Z_p}S_p(E_{\rm c},e)=0 .
 \tag{EW.6}
\]
This tensor identity is literal: a $\mathbb Z_p$-basis of the
finite free module $O_p$ identifies its coefficient and continuous
cochains with finitely many copies of the untwisted ones, globally
and locally, and therefore identifies their kernels.
Corestriction composed with restriction is multiplication by
$[E_{\rm c}:K]=3$ on the original cohomology.
Three is a unit of $O_p$, so restriction is injective.
This proves the three residue-zero rows in $Z_{\rm cub}$.
Equivalently one can apply the integral $C_3$ Fourier/Shapiro
decomposition; its only denominator is the same unit three.
No Fourier decomposition at coefficient three is invoked here.
The source's separate exclusion of this residue by \texttt{GD.gen}
does not exclude it from (TC.1), whose weight range is all $1\leq e<T$.

\proofstep{4. The cubic and trivial rows at residue $267$.}
AC-G103, equation (G103.12), already proves the exact residual
strict vanishing
\[
             \mathcal S(K,k(\bar c_p\bar\chi^{\pm1}))=0.
 \tag{EW.7}
\]
Its definition (G103.2) uses the full local zero condition at $p$
and inertia zero at every other finite place.
It holds over every finite field containing the displayed values.
Its proof uses only the already certified fixed cubic class group
and its exact unit localization, without an endpoint hypothesis.

For the trivial cubic label, the analogous residual kernel
$\mathcal S(K,k(\bar c_p))$ is zero as well.
Indeed Kummer theory identifies a strict class with an element
whose valuations are divisible by $p$ at every finite prime.
The all-primes valuation sequence has unit term
$\mathcal O_K^\times/\mathcal O_K^{\times p}$ and class term
$\operatorname{Cl}(K)[p]$.  The norm-Euclidean ring
$\mathbb Z[\omega]$ has class number one and units $\mu_6$,
so both terms are zero at $p=269$.
These are the elementary base-field facts already used in AC-G22.
Extending scalars from $\mathbb F_p$ to the actual finite residue
field preserves this kernel, by the same finite-free cochain argument.

For $e\equiv267\bmod268$, the reduction of
$D=(E_p/O_p)(\chi^j c_p^{-e})$ is therefore strict-zero for
each $j=0,1,2$.
We give the integral lift rather than change the fixed character
in an earlier theorem silently.
At $p$, for $j=1,2$ the cubic residual inertia has order three
and cannot cancel the order-dividing-$268$ cyclotomic factor.
For $j=0$, $\bar c_p$ is nontrivial on $G_{K_p}$, since $K_p$
is unramified quadratic and the $p$-cyclotomic extension is
totally ramified.  In both cases $D^{G_{K_p}}=0$.
Away from $p$, the inertia character has order dividing three,
so its invariants are zero or all of the divisible module $D$.
For an original uniformizer $\varpi$ the local maps
$H^1(J,D[\varpi])\to H^1(J,D)$ are consequently injective.
The long exact sequence lifts every strict $\varpi$-torsion class
to a residual strict class, which is zero.
Each continuous cocycle has finite image in $E_p/O_p$, hence is
killed by a power of $\varpi$.
There can be no nonzero strict class.
This proves the three residue-$267$ rows.
The two positive-weight progressions in $Z_{\rm cub}$ are disjoint,
with counts $\lfloor(T-1)/268\rfloor$ and $\lfloor T/268\rfloor$.

\proofstep{5. The old coefficient-three trivial row, including ramified $O_3$.}
AC-K11 explicitly proves
\[
                         S_3(K,e)=0\qquad(e\geq1).
 \tag{EW.8}
\]
Its proof retains the cited regular-prime higher-K-theory statement
and the injection in (K10.1); neither an unproved all-prime
regularity claim nor a coefficient-field vanishing is substituted.
Its strict group has coefficient $\mathbb Q_3/\mathbb Z_3(-e)$,
unramified outside three and zero in full local cohomology there.
The ring $O_3$ in (FC.1) can be ramified.
It is nevertheless finite free of rank $b_3$ over $\mathbb Z_3$.
The direct-sum continuous-cochain argument in Step3 proves
\[
 \mathcal S_{3,1}(e)
       =O_3\otimes_{\mathbb Z_3}S_3(K,e)=0 ,
 \tag{EW.9}
\]
after the exact added-prime comparison in Step1.
This uses no division by $|G_u|$, which is divisible by three.
It proves the final row and its count $T-1$.

\proofstep{6. A single absolute error and the unchanged common window.}
For every removed label, (EU.2) gives
\[
 \frac{\log\#k_{O_\ell}}{b_\ell}\max(0,i_{\ell,\eta,e})
 \leq\frac{\log\#H^2(C_{\ell,\eta,e})}{b_\ell}
 =\frac{\left|\log\#H^2(C_{\ell,\eta,e})
                  -\log\#\mathcal S_{\ell,\eta}(e)\right|}{b_\ell}.
 \tag{EW.10}
\]
All terms are nonnegative, so the union is one subsum of the
full absolute comparison (FL.4), including both coefficient
primes and every original normalization.
Summing that bound once over $1\leq e<T$ gives (EW.3) by
$\sum e\leq T^2/2$ and $\log(8eD)\leq\log(8TD)$.
The host estimates in AC-G143/AC-G145 give the stated exponent.
No positive index is individually set to zero.

Given the new common-window premise with slack $s>0$, choose
$U_0$ so that (EW.3) is at most
$(s/2)\exp(u^\varepsilon/2)$ for $u\geq U_0$.
Increase $X_0$ to $\max(X_0,U_0^J)$.
The exact same $J$, page, pair, original coefficient rings,
host and common $W$ then satisfy the full old premise with
slack $s/2$.  Simultaneous finite labels use the maximum of
their uniform error constants, not intersections of independently
chosen windows.  Nonnegativity gives every reverse implication.
All unremoved terms and the mixed input $\mathcal R_{\rm cub}$ remain.
\end{proof}

\paragraph{Audit of the other established lift families.}
This paragraph concerns the audited original-character families,
not a claimed classification of all possible strict-zero characters.
AC-G126 has the tame factor $\chi_{-20}$, and AC-G130/AC-G134
have the primitive quartic factor modulo five.
AC-G141 retains that quartic factor while multiplying it by $\chi_8$.
Every one is nontrivial on inertia at five.
The actual (FG.3) hosts are unramified at five: their prescribed
primes satisfy $q,d>7$.
Their extra Teichmuller and principal-unit factors are unramified
there, so these families supply no character label of that host.

For the audited conductor-four and conductor-sixteen families,
suppose an entire lifted finite character factors through $G_u$.
The image of $G_u$ has order $144\cdot2^{h(u)}$, prime to $269$.
All fixed finite and Teichmuller factors have order prime to $269$.
Consequently the principal-unit factor, if the product is finite,
is both of prime-to-$269$ order and of $269$-power order; it is trivial.
The native prime congruences are $q\equiv3\bmod8$ and $d\equiv7\bmod16$.
Since $269\equiv5\bmod8$ and $269\equiv13\bmod16$, neither is $269$.
Thus the noncubic factor of $H_u/K$ is unramified at $269$,
and its inertia image has only the cubic order three.
The cyclotomic image of inertia on $G_K$ is $\mathbb F_{269}^\times$.
Restricting to the kernel of the cubic character does not shrink
this image, since three and $268$ are coprime.
A Teichmuller factor that descends to the host must therefore have
exponent divisible by $268$.
There are no further actual labels from those additional lifts.

AC-G22's coefficient-three vanishing occurs at weight zero,
outside (TC.1); its positive-weight cubic group is not asserted zero.
The rank-two coefficient-field statements AC-K29 and the adjoint
rigidity statements do not assert integral strict vanishing for
the rank-one modules (FC.2).
The other class-sector bounds and host model-length bounds are
not treated as zero groups.
No exclusion of arbitrary new factors at the variable primes $q,d$
is asserted.  They remain in the arithmetic gate.

\paragraph{Source scope.}
The exact new transports above use the complete existing AC-K31
fixed-field proof, AC-K11 and its stated higher-K input,
the actual residual statement and unit proof (G103.2), (G103.12),
and the original-coefficient theorems AC-G114/115/125.
AC-G142/144 and AC-G145 supply the conductor-sixteen rows.
AC-K25, (FL.4), (EU.2), and the literal source \texttt{WinPair}
conditions retain their earlier applied-source scope.
The finite-cover dictionary remains the retained SGA~1,
Expos\'e~V, \S\S4 and 7 dictionary; Kummer descent is the
retained Stacks tag~03PK input.
No new primary-source proof attribution, finite number-field
calculation, or Lean execution is claimed.
